\scanpage{72}
% lines 1–5, diagonally ruled attempt and right-hand continuation
\deleted{Faisons $\unsure{\to}$}
\[
\deleted{x_1=\frac{1}{\unsure{\chi|\gamma}_1}T_1\qquad x_2=\frac{1}{\chi}T_2\qquad(!)}
\]
\deleted{on trouve}
\[
\deleted{T=\chi\underbrace{\bigl(\tfrac1\chi T_1\bigr)*\bigl(\tfrac1\chi T_2\bigr)}_{\frac1\chi T_1*T_2}-}
\]
\margin{Faisons $x_{1}=T_1*T_2$, d'où}
% lines 6–8, cancellation
d'où
\[
\chi^2T=\chi\,T_1*T_2
\]
d'où en divisant par $\chi$
\[
\boxed{\chi T=T_1*T_2}
\]
% lines 9–13, induced application
\uline{NB} L'application
\[
\begin{aligned}
E_1\times E_2&\mathbin{-}E_1\otimes E_2\longrightarrow E_1*E_2\\
x_1,x_2&\longmapsto x_1\otimes x_2,\ \ x_1\otimes x_2\longmapsto x_1*x_2
\end{aligned}
\]
applique
\[
P_1\times P_2\longrightarrow P=\psi^{-1}(1)
\]
% lines 14–19, elements and product expansion
On aura, si
\[
\begin{aligned}
x_1&\in\Gamma P_1,\ u_1\in\Gamma L_1\\
x_2&\in\Gamma P_2,\ u_2\in\Gamma L_2
\end{aligned}
\]
\[
\begin{aligned}
(\underset{\substack{\rotatebox{90}{$\in$}\\\deleted{\unsure{\Gamma P_1}}}}{x_1}
+\underset{\substack{\rotatebox{90}{$\in$}\\\Gamma L_1}}{u_1})*
(\underset{\substack{\rotatebox{90}{$\in$}\\\Gamma P_2}}{x_2}
+\underset{\substack{\rotatebox{90}{$\in$}\\\Gamma L_2}}{u_2})
&=\underbrace{x_1*x_2}_{\Gamma P}
+\underbrace{u_1*x_2+x_1*u_2}_{\Gamma(\deleted{E}L_1\otimes L_2)}
+\underbrace{u_1*u_2}_{\shortparallel\,\chi u_1\otimes u_2}\\
&=x_1*x_2+\bigl[u_1\otimes(\overbrace{\chi x_2-\deleted{\unsure{\psi_2(x_2)|\psi(x_2)}}T_2}^{-b_2})
+(\overbrace{\chi x_1-\deleted{\unsure{\psi_1(x_1)|\psi(x_1)}}T_1}^{-b_1})\otimes u_2\\
&\hspace{9em}+\chi\,u_1\otimes u_2\bigr]
\end{aligned}
\]
% lines 20–23, lower left coordinates, boxed calculation and right-edge note
\margin{\uline{NB} $x_1=(0,1)\quad T_1=(b_1,\chi)$\\ $\chi x_1=(0,\chi)$\\ \deleted{$\chi x_1-T_1$}\\ $\chi x_1-T_1=(-b_1,0)$}
\[
\boxed{\begin{aligned}
(x_1+u_1)*(x_2+u_2)&=x_1*x_2+(-u_1\otimes b_2-b_1\otimes u_2+\chi\,u_1\otimes u_2)\\
&=x_1*x_2+(-u_1\otimes b_2\mathbin{\unsure{-|+}}b_1\otimes u_2+\chi\,u_1\otimes u_2)
\end{aligned}}
\]
\margin{\unsure{ces signes vont s'éliminer si} \unsure{subst.}}
% lines 24–27, slanted lower-left note and baseline
% slanted
\margin{\uline{NB} Si $\{b_1\}$ est une base de $L_1$, alors $x_2\longrightarrow x_1*x_2$ est un iso $P_2\simeq P_1*P_2$ de torseurs, compatible \unsure{à|avec} l'isom $u_2\longmapsto -b_1\otimes u_2$ de $L_2$ sur $L_1\otimes L_2\simeq O\otimes L_2=L_2$}
$b_1,b_2$ sont les éléments de $\deleted{E}L_1,L_2$ qui définissent $T_1,T_2$ en termes des $x_1,x_2$
\[
T_1=\chi x_1+b_1,\quad T_2=\chi x_2+b_2
\]
% archive sheet number: 71
